\documentclass[10pt,aspectratio=169]{beamer}
\input{../../_en_preambulo_beamer}
% Comandos y macros estadísticas compartidas para las presentaciones Beamer en inglés (Metropolis)

% Conjuntos numéricos básicos
\newcommand{\R}{\mathbb{R}}
\newcommand{\N}{\mathbb{N}}
\newcommand{\Q}{\mathbb{Q}}
\newcommand{\Z}{\mathbb{Z}}
\newcommand{\set}[1]{\left\{ #1 \right\}}
\newcommand{\abs}[1]{\left|#1\right|}
\newcommand{\norm}[1]{\left\| #1 \right\|}

% Comandos estadísticos y probabilísticos del curso
\newcommand{\Var}{\operatorname{Var}}
\newcommand{\cov}{\operatorname{Cov}}
\newcommand{\corr}{\rho}
\newcommand{\comb}[2]{\begin{pmatrix} #1 \\ #2 \end{pmatrix}}
\newcommand{\s}{\sigma}
\newcommand{\particion}{\mathfrak{P}}
\newcommand{\card}[1]{\#\left( #1 \right)}
\newcommand{\seq}[3]{\set{#1}_{#2}^{#3}}
\newcommand{\E}{\mathbb{E}}
\newcommand{\Prob}{\mathbb{P}}

% Entornos pedagógicos y cajas en inglés académico natural (soportando tanto nombres en inglés como en español)
\newenvironment{discussion}{%
  \begin{alertblock}{Discussion Question}%
}{%
  \end{alertblock}%
}
\newenvironment{discusion}{%
  \begin{alertblock}{Discussion Question}%
}{%
  \end{alertblock}%
}

\newenvironment{reflection}{%
  \begin{alertblock}{Classroom Reflection}%
}{%
  \end{alertblock}%
}
\newenvironment{reflexion}{%
  \begin{alertblock}{Classroom Reflection}%
}{%
  \end{alertblock}%
}

\newenvironment{paradox}{%
  \begin{alertblock}{Exploratory Paradox}%
}{%
  \end{alertblock}%
}
\newenvironment{paradoja}{%
  \begin{alertblock}{Exploratory Paradox}%
}{%
  \end{alertblock}%
}

\newenvironment{motivation}{%
  \begin{exampleblock}{Motivation: Data Science in Practice}%
}{%
  \end{exampleblock}%
}
\newenvironment{motivacion}{%
  \begin{exampleblock}{Motivation: Data Science in Practice}%
}{%
  \end{exampleblock}%
}

\newenvironment{quickchallenge}{%
  \begin{block}{Quick Team Challenge}%
}{%
  \end{block}%
}
\newenvironment{retoclase}{%
  \begin{block}{Quick Team Challenge}%
}{%
  \end{block}%
}


\title[Fundamental Concepts]{Section 5.8: Fundamental Concepts of Inferential Statistics}
\subtitle{Estimators, standard error, and sampling limits}
\author[J. Castillo Colmenares]{Juliho Castillo Colmenares}
\institute{Tecnológico de Monterrey}
\date{}

\begin{document}
\begin{frame}[plain]\vspace{-0.8cm}\titlepage\end{frame}

\begin{frame}{Roadmap --- Chapter 5}
  \footnotesize
  \begin{itemize}\itemsep=0.08cm
    \item \textbf{05.00} Inferential introduction
    \item \textbf{05.01--05.05} Sampling distributions
    \item \textbf{05.08} Fundamental concepts \emph{(today)}
    \item \textbf{05.09} Z and t statistics
  \end{itemize}
  \begin{block}{Learning objective}
    Connect estimators, sampling distributions, and standard error with the law
    of large numbers and the central limit theorem.
  \end{block}
\end{frame}

\begin{frame}{Motivation: estimating a parameter}
  \footnotesize
  An estimator is a statistic used to approximate an unknown population
  parameter. We write $\hat\theta$ for the parameter $\theta$.
  \begin{alertblock}{Guiding question}
    Is the estimator centered on the true value, and how variable is it?
  \end{alertblock}
\end{frame}

\begin{frame}{Unbiased estimator}
  \footnotesize
  \begin{block}{Definition}
    $\hat\theta$ is unbiased for $\theta$ if
    \[
      E(\hat\theta)=\theta.
    \]
  \end{block}
  \begin{exampleblock}{Reading example}
    The sample mean satisfies $E(\bar X)=\mu$.
  \end{exampleblock}
\end{frame}

\begin{frame}{Sampling distribution}
  \footnotesize
  The sampling distribution of a statistic is its probability distribution when
  computed over all possible samples of size $n$.
  \begin{block}{Object of study}
    It describes how the statistic changes when sampling is repeated, not the
    variation of observations inside one sample.
  \end{block}
\end{frame}

\begin{frame}{Standard error}
  \footnotesize
  Standard error is the standard deviation of a statistic's sampling
  distribution:
  \[
    SE(\bar X)=\frac{\sigma}{\sqrt n},
    \qquad SE(\bar X)\approx\frac{s}{\sqrt n}.
  \]
  \begin{alertblock}{Reading}
    It measures estimation precision, not individual-observation spread.
  \end{alertblock}
\end{frame}

\begin{frame}{Law of large numbers}
  \footnotesize
  For iid variables with mean $\mu$,
  \[
    \lim_{n\to\infty}P(|\bar X_n-\mu|<\varepsilon)=1.
  \]
  With large samples, the sample mean approaches the population mean in
  probability.
\end{frame}

\begin{frame}{Central limit theorem}
  \footnotesize
  If the iid variables have mean $\mu$ and finite variance $\sigma^2$,
  \[
    \frac{\bar X-\mu}{\sigma/\sqrt n}\xrightarrow{d}N(0,1).
  \]
  In practice the CLT often works from $n\ge30$, depending on population shape.
\end{frame}

\begin{frame}{Importance for inference}
  \footnotesize
  The CLT allows us to:
  \begin{itemize}\itemsep=0.1cm
    \item use the normal distribution for means from non-normal populations;
    \item justify Z and t tests;
    \item construct confidence intervals;
    \item understand approximately normal patterns in nature.
  \end{itemize}
\end{frame}

\begin{frame}{Computational bridge 1/4: estimator}
  \footnotesize
  \[
    \bar X=\frac1n\sum_{i=1}^nX_i,
    \qquad E(\bar X)=\mu.
  \]
  The expected value centers the estimator at the parameter of interest.
\end{frame}

\begin{frame}{Computational bridge 2/4: variability}
  \footnotesize
  \[
    \operatorname{Var}(\bar X)=\frac{\sigma^2}{n},
    \qquad SE(\bar X)=\frac{\sigma}{\sqrt n}.
  \]
  Increasing $n$ reduces the variability of the sample mean.
\end{frame}

\begin{frame}{Computational bridge 3/4: standardization}
  \footnotesize
  \[
    Z_n=\frac{\bar X-\mu}{\sigma/\sqrt n}.
  \]
  \begin{alertblock}{Limit}
    The CLT states that $Z_n$ approaches a standard normal distribution, making
    known probabilities and quantiles available.
  \end{alertblock}
\end{frame}

\begin{frame}{Computational bridge 4/4: sample size}
  \footnotesize
  \begin{tabular}{lll}
    \hline
    Change & Effect on $SE(\bar X)$ & Reading \\
    \hline
    $n\to4n$ & divided by $2$ & higher precision \\
    $n\to n/4$ & multiplied by $2$ & lower precision \\
    \hline
  \end{tabular}
  \begin{block}{Conclusion}
    Precision improves with $n$, but non-sampling errors remain.
  \end{block}
\end{frame}

\begin{frame}{Solved example 1/4 --- unbiased mean}
  \footnotesize
  The reading considers $\bar X=(1/n)\sum X_i$ as an estimator of $\mu$.
  \begin{block}{Statement}
    Compute $E(\bar X)$ using linearity of expectation.
  \end{block}
\end{frame}

\begin{frame}{Solved example 1/4 --- solution}
  \footnotesize
  \[
    E(\bar X)=E\!\left(\frac1n\sum X_i\right)
    =\frac1n\sum E(X_i)=\frac1n\,n\mu=\mu.
  \]
  \begin{alertblock}{Result}
    The sample mean is an unbiased estimator of the population mean.
  \end{alertblock}
\end{frame}

\begin{frame}{Solved example 2/4 --- sampling distribution}
  \footnotesize
  For iid variables with mean $\mu$ and variance $\sigma^2$, the reading studies
  the distribution of $\bar X$.
  \begin{block}{Statement}
    Determine its mean and variance.
  \end{block}
\end{frame}

\begin{frame}{Solved example 2/4 --- solution}
  \footnotesize
  \[
    E(\bar X)=\mu,\qquad \operatorname{Var}(\bar X)=\frac{\sigma^2}{n}.
  \]
  \begin{exampleblock}{Interpretation}
    The sample mean becomes less variable as the sample size grows.
  \end{exampleblock}
\end{frame}

\begin{frame}{Solved example 3/4 --- exponential waiting times}
  \footnotesize
  A waiting time is exponential with mean $\mu=5$ minutes. For $n=50$, the
  reading applies the CLT to the sample mean.
  \begin{block}{Statement}
    Find the approximate distribution of $\bar X$.
  \end{block}
\end{frame}

\begin{frame}{Solved example 3/4 --- solution}
  \footnotesize
  \[
    \bar X\approx N\!\left(5,\frac{5^2}{50}\right)=N(5,0.5).
  \]
  \begin{alertblock}{Interpretation}
    Although the population is exponential, the mean of 50 observations is
    approximately normal.
  \end{alertblock}
\end{frame}

\begin{frame}{Reading application 4/4 --- solution}
  \footnotesize
  \[
    \text{sample}\to\text{statistic}\to\text{sampling distribution}
    \to SE\to\text{normal approximation}\to\text{inference}.
  \]
  This sequence reproduces the conceptual structure of the reading.
\end{frame}

\begin{frame}{Synthesis of the section}
  \footnotesize
  \begin{itemize}\itemsep=0.12cm
    \item Unbiasedness centers the estimator.
    \item Standard error quantifies precision.
    \item The law of large numbers explains convergence.
    \item The CLT supplies a normal limit.
  \end{itemize}
\end{frame}

\begin{frame}{Curricular transition}
  \footnotesize
  \begin{block}{Established}
    We can describe how an estimator behaves when sampling is repeated.
  \end{block}
  \begin{alertblock}{Next section}
    \textbf{05.09 Z and t Statistics}: standardization for tests about means
    when the standard deviation is known or unknown.
  \end{alertblock}
\end{frame}
\end{document}
